\documentclass[12pt]{article}
\usepackage{amsmath,amsfonts,amssymb}
\usepackage{geometry}
\usepackage{mathtools}
\usepackage{enumitem}
\usepackage{hyperref}
\usepackage{amsthm}

\geometry{margin=1in}
\setlength{\parskip}{1em}
\setlength{\parindent}{0pt}

\title{The Conjugate Gradient Method in the Most General Setting}
\author{}
\date{}

\begin{document}

\maketitle

\section{Problem Setting}

Let \( \mathcal{H} \) be a real Hilbert space with inner product \( \langle \cdot, \cdot \rangle \) and norm \( \| \cdot \| \). Let \( A : \mathcal{H} \to \mathcal{H} \) be a bounded, self-adjoint, and positive-definite linear operator, and let \( f \in \mathcal{H} \).

We consider the problem of finding \( u \in \mathcal{H} \) such that:
\[
A u = f.
\]

\section{Variational Formulation}

This problem is equivalent to minimizing the quadratic functional:
\[
J(v) = \frac{1}{2} \langle A v, v \rangle - \langle f, v \rangle.
\]

Because \( A \) is self-adjoint and positive-definite, \( J \) is strictly convex and has a unique minimizer, which is the solution of \( A u = f \).

\section{The Conjugate Gradient Method}

Let \( u_0 \in \mathcal{H} \) be an initial guess. We define:

\begin{align*}
r_0 &= f - A u_0 \quad \text{(residual)}, \\
p_0 &= r_0 \quad \text{(initial search direction)}.
\end{align*}

Then, for \( k = 0, 1, 2, \dots \), iterate:

\begin{enumerate}[label=(\arabic*)]
    \item Compute step size:
    \[
    \alpha_k = \frac{\langle r_k, r_k \rangle}{\langle A p_k, p_k \rangle}.
    \]
    
    \item Update solution:
    \[
    u_{k+1} = u_k + \alpha_k p_k.
    \]
    
    \item Update residual:
    \[
    r_{k+1} = r_k - \alpha_k A p_k.
    \]
    
    \item Check for convergence: if \( \|r_{k+1}\| \) is below a given tolerance, stop.
    
    \item Compute conjugation factor:
    \[
    \beta_k = \frac{\langle r_{k+1}, r_{k+1} \rangle}{\langle r_k, r_k \rangle}.
    \]
    
    \item Update direction:
    \[
    p_{k+1} = r_{k+1} + \beta_k p_k.
    \]
\end{enumerate}

\section{Remarks}

\begin{itemize}
    \item The method only requires the application of \( A \) to elements in \( \mathcal{H} \), not an explicit representation.
    
    \item Convergence is guaranteed in at most \( n \) steps if \( \mathcal{H} \) is \( n \)-dimensional, but in infinite dimensions, the method must be truncated based on a tolerance.
    
    \item In practice, for problems in infinite-dimensional spaces (e.g., PDEs), one typically discretizes the problem using finite-dimensional subspaces (e.g., finite element or spectral methods) and applies CG to the resulting system.
    
    \item If \( A \) is not positive-definite or not self-adjoint, the conjugate gradient method is not applicable. Alternatives include MINRES or GMRES.
\end{itemize}

\section{Discretization}

To compute numerically, we select a sequence of finite-dimensional subspaces \( V_h \subset \mathcal{H} \) and solve the projected problem:
\[
\text{Find } u_h \in V_h \text{ such that } \langle A u_h, v_h \rangle = \langle f, v_h \rangle \quad \forall v_h \in V_h.
\]
This leads to a linear system \( A_h \mathbf{u}_h = \mathbf{f}_h \), where CG can be applied in matrix or operator form.

\section{Derivation}

\subsection{Motivation: Steepest Descent}

Starting from an initial guess \( u_0 \), one might attempt to minimize \( J \) using steepest descent. The direction of steepest descent is given by the negative gradient of \( J \) at \( u_k \), which is
\[
\nabla J(u_k) = A u_k - f = -r_k,
\]
where we define the residual
\[
r_k = f - A u_k.
\]
Thus, the steepest descent update is
\[
u_{k+1} = u_k + \alpha_k r_k,
\]
where \( \alpha_k \) is chosen to minimize \( J(u_k + \alpha r_k) \). Setting the derivative of \( J(u_k + \alpha r_k) \) with respect to \( \alpha \) to zero, we get:
\[
\frac{d}{d\alpha} J(u_k + \alpha r_k) = \langle A(u_k + \alpha r_k), r_k \rangle - \langle f, r_k \rangle = \langle A r_k, r_k \rangle \alpha - \langle r_k, r_k \rangle.
\]
Setting this to zero yields:
\[
\alpha_k = \frac{\langle r_k, r_k \rangle}{\langle A r_k, r_k \rangle}.
\]

\textbf{Problem:} Steepest descent suffers from slow convergence, especially when \( A \) is ill-conditioned. Hence, we improve the method by enforcing conjugacy of directions.

\subsection{Conjugate Directions and Optimality}

We seek to build a sequence of directions \( \{p_k\} \subset \mathcal{H} \) such that:

\begin{itemize}
    \item \( p_k \in \operatorname{span}\{r_0, A r_0, \dots, A^{k} r_0\} \) (Krylov subspace),
    \item \( \langle A p_i, p_j \rangle = 0 \) for \( i \neq j \) (mutual \( A \)-conjugacy).
\end{itemize}

Then, the iterates are defined by:
\[
u_{k+1} = u_k + \alpha_k p_k,
\]
and we choose \( \alpha_k \) to minimize \( J(u_k + \alpha p_k) \). This gives:
\[
\alpha_k = \frac{\langle r_k, p_k \rangle}{\langle A p_k, p_k \rangle}.
\]
We now define the new residual:
\[
r_{k+1} = r_k - \alpha_k A p_k.
\]

\subsubsection*{Why this residual update?}

Note:
\[
r_{k+1} = f - A u_{k+1} = f - A(u_k + \alpha_k p_k) = r_k - \alpha_k A p_k.
\]
So this is not a guess — it's an exact consequence of the update to \( u_{k+1} \).

\subsection{Constructing Conjugate Directions}

We now derive how to build \( p_{k+1} \) from \( r_{k+1} \) and \( p_k \) so that \( p_{k+1} \) remains \( A \)-conjugate to the previous directions.

We propose:
\[
p_{k+1} = r_{k+1} + \beta_k p_k,
\]
and we determine \( \beta_k \) such that \( \langle A p_{k+1}, p_k \rangle = 0 \). Substituting:
\[
\langle A (r_{k+1} + \beta_k p_k), p_k \rangle = 0 \Rightarrow \langle A r_{k+1}, p_k \rangle + \beta_k \langle A p_k, p_k \rangle = 0.
\]
Solving for \( \beta_k \), we get:
\[
\beta_k = - \frac{\langle A r_{k+1}, p_k \rangle}{\langle A p_k, p_k \rangle}.
\]

\textbf{Alternative formula:} In practice, we use a numerically cheaper and equivalent expression (which follows from orthogonality of residuals):
\[
\beta_k = \frac{\langle r_{k+1}, r_{k+1} \rangle}{\langle r_k, r_k \rangle}.
\]

This is valid due to:
\[
\langle r_{k+1}, p_k \rangle = 0,
\]
which can be shown inductively.

\subsection{Final Conjugate Gradient Algorithm}

We summarize the derived steps:

\begin{enumerate}[label=(\arabic*)]
    \item Choose \( u_0 \in \mathcal{H} \), compute \( r_0 = f - A u_0 \), set \( p_0 = r_0 \).
    
    \item For \( k = 0, 1, 2, \dots \), do:
    \begin{align*}
    \alpha_k &= \frac{\langle r_k, r_k \rangle}{\langle A p_k, p_k \rangle}, \\
    u_{k+1} &= u_k + \alpha_k p_k, \\
    r_{k+1} &= r_k - \alpha_k A p_k, \\
    \beta_k &= \frac{\langle r_{k+1}, r_{k+1} \rangle}{\langle r_k, r_k \rangle}, \\
    p_{k+1} &= r_{k+1} + \beta_k p_k.
    \end{align*}
    
    \item Stop when \( \| r_k \| \) is below tolerance.
\end{enumerate}

\subsection{Geometric and Optimality Interpretation}

At each iteration, \( u_k \) minimizes \( J \) over the \( k \)-th Krylov subspace:
\[
\mathcal{K}_k = \operatorname{span}\{r_0, A r_0, \dots, A^{k-1} r_0\}.
\]

Thus, CG is a Galerkin method in a Krylov subspace, selecting the best approximation in that space with respect to the energy norm \( \| v \|_A := \sqrt{\langle A v, v \rangle} \).


\section{Further details on deriving \(\alpha_k\)}

We aim to minimize the energy functional
\[
J(v) = \frac{1}{2} \langle A v, v \rangle - \langle f, v \rangle
\]
along the search direction \( d_k \in \mathcal{H} \), starting from the current iterate \( u_k \). That is, we define:
\[
\phi(\alpha) := J(u_k + \alpha d_k)
\]
and we want to choose \( \alpha_k \) that minimizes \( \phi(\alpha) \).

\paragraph{Step 1: Expand \( \phi(\alpha) \).}

Using bilinearity and symmetry of the inner product and of \( A \), we have:
\begin{align*}
\phi(\alpha)
&= \frac{1}{2} \langle A(u_k + \alpha d_k), u_k + \alpha d_k \rangle - \langle f, u_k + \alpha d_k \rangle \\
&= \frac{1}{2} \left[ \langle A u_k, u_k \rangle + 2\alpha \langle A u_k, d_k \rangle + \alpha^2 \langle A d_k, d_k \rangle \right] \\
&\quad - \langle f, u_k \rangle - \alpha \langle f, d_k \rangle.
\end{align*}

\paragraph{Step 2: Differentiate with respect to \( \alpha \).}

\[
\phi'(\alpha) = \langle A u_k, d_k \rangle - \langle f, d_k \rangle + \alpha \langle A d_k, d_k \rangle.
\]

\paragraph{Step 3: Specialize to \( d_k = r_k = f - A u_k \).}

Substitute \( d_k = r_k \), and note that \( A u_k = f - r_k \), so:
\[
\langle A u_k, r_k \rangle = \langle f, r_k \rangle - \langle r_k, r_k \rangle.
\]

Then,
\begin{align*}
\phi'(\alpha)
&= \left( \langle f, r_k \rangle - \langle r_k, r_k \rangle \right) - \langle f, r_k \rangle + \alpha \langle A r_k, r_k \rangle \\
&= - \langle r_k, r_k \rangle + \alpha \langle A r_k, r_k \rangle.
\end{align*}

\paragraph{Step 4: Set derivative to zero.}

Set \( \phi'(\alpha_k) = 0 \) to minimize \( \phi \). Then:
\[
\alpha_k = \frac{\langle r_k, r_k \rangle}{\langle A r_k, r_k \rangle}.
\]

\section*{Further derivation of \texorpdfstring{$\alpha_k$}{alpha\_k} in the Conjugate Gradient Method}

At iteration \( k \), we seek to compute the next approximation \( u_{k+1} \) as:
\[
u_{k+1} = u_k + \alpha_k p_k,
\]
where \( p_k \in \mathcal{H} \) is the current search direction. To find the optimal step size \( \alpha_k \), we minimize the energy functional along this direction:
\[
\phi(\alpha) := J(u_k + \alpha p_k),
\]
where
\[
J(v) = \frac{1}{2} \langle A v, v \rangle - \langle f, v \rangle.
\]

\paragraph{Step 1: Expand \( \phi(\alpha) \).}

\begin{align*}
\phi(\alpha)
&= \frac{1}{2} \langle A(u_k + \alpha p_k), u_k + \alpha p_k \rangle - \langle f, u_k + \alpha p_k \rangle \\
&= \frac{1}{2} \left[ \langle A u_k, u_k \rangle + 2\alpha \langle A u_k, p_k \rangle + \alpha^2 \langle A p_k, p_k \rangle \right] \\
&\quad - \langle f, u_k \rangle - \alpha \langle f, p_k \rangle.
\end{align*}

\paragraph{Step 2: Differentiate with respect to \( \alpha \).}

\[
\phi'(\alpha) = \langle A u_k, p_k \rangle - \langle f, p_k \rangle + \alpha \langle A p_k, p_k \rangle.
\]

\paragraph{Step 3: Use the definition of the residual.}

Recall the residual at step \( k \) is defined as:
\[
r_k = f - A u_k \quad \Rightarrow \quad \langle A u_k, p_k \rangle = \langle f, p_k \rangle - \langle r_k, p_k \rangle.
\]
Substituting into \( \phi'(\alpha) \), we get:
\begin{align*}
\phi'(\alpha)
&= (\langle f, p_k \rangle - \langle r_k, p_k \rangle) - \langle f, p_k \rangle + \alpha \langle A p_k, p_k \rangle \\
&= - \langle r_k, p_k \rangle + \alpha \langle A p_k, p_k \rangle.
\end{align*}

\paragraph{Step 4: Solve for the minimizer.}

Set \( \phi'(\alpha_k) = 0 \), then:
\[
\alpha_k = \frac{\langle r_k, p_k \rangle}{\langle A p_k, p_k \rangle}.
\]

\paragraph{Residual Update.}

Once \( \alpha_k \) is computed, the new iterate is:
\[
u_{k+1} = u_k + \alpha_k p_k,
\]
and the new residual becomes:
\begin{align*}
r_{k+1} &= f - A u_{k+1} = f - A(u_k + \alpha_k p_k) \\
&= f - A u_k - \alpha_k A p_k = r_k - \alpha_k A p_k.
\end{align*}

\paragraph{Conclusion:}

The step size \( \alpha_k \) is chosen to minimize the energy functional along the \( A \)-conjugate direction \( p_k \), and the residual update follows directly from the definition of \( u_{k+1} \). This ensures that \( r_{k+1} \perp p_k \), preserving the orthogonality structure that underpins the efficiency of the conjugate gradient method.



\end{document}
